\section{Introduction}

Safety guards are classifiers placed before or after a language model to detect
harmful content, jailbreaks, and prompt injections. Their deployment objective is
not merely high accuracy on one benchmark: a useful guard must preserve benign
traffic, catch harmful behavior, and transfer across datasets whose annotation
policies and lexical cues differ. This makes the training objective consequential,
but also makes objective comparisons unusually easy to confound.

Recent guard pipelines use DPO to emphasize hard examples or preference-unstable
rollouts \citep{liu2025guardreasoner,ferrao2026aims,liu2026dtguard}. These results
establish that preference optimization can participate in effective guard
training. They do not isolate what DPO's frozen reference contributes when the
``preference'' is mechanically derived from a binary class label. In that case,
chosen is the correct verdict token and rejected is the incorrect token; no new
human or policy preference has been added.

We ask a narrower question: \emph{what does reference centering itself buy a
binary safety guard?} We compare three objectives from the same Stage-1 adapter:
full-vocabulary verdict cross-entropy (VerdictCE), a temperature-matched pairwise
margin loss (PairCE), and DPO. PairCE and DPO see identical examples and share the
same inverse temperature, so their difference is the subtraction of the frozen
Stage-1 margin. A second factorial axis compares high-entropy examples with a
disjoint source/label-matched random control.

Asking that question carefully turns out to be harder than running it. Once the
reduction is in hand, DPO and PairCE differ by a single per-example gradient
weight, and we show that at initialization the ratio of those weights is
near-constant on a real guard panel: mean $1.30$ with a coefficient of variation
of $0.17$ (\Cref{lem:rescale}). Reference centering therefore \emph{begins} as a
larger effective learning rate, not as a differential re-emphasis. Any comparison
at a fixed step count silently mixes ``better objective'' with ``faster
training.'' Worse, the natural estimand is underpowered: the measured
within-checkpoint seed standard deviation on the parent suite is $0.036$, and with
four checkpoints as the independent unit a one-sided bound needs an effect above
$0.035$ --- larger than an explicit KL anchor achieves on the same data. A study
built on that estimand could not tell a null from a real small effect.

We make the confound the hypothesis. A pure learning-rate rescale moves a run
along its own represented/transfer trade-off curve without moving the curve, so we
compare curves at matched represented ranking instead of points at matched step
count. Each cell contributes a trajectory over a saved checkpoint ladder, and the
primary estimand is the transfer gap between objectives at equal represented AP.
It is exactly zero under a rescale and nonzero only when one objective buys
transfer at equal represented ranking.

The planned contribution is fivefold. First, we derive the exact one-token DPO
reduction and state its interpretation boundary. Second, we prove the
near-uniform-rescale lemma and draw its two design consequences. Third, we specify
a matched $3\times2$ objective-by-selection design over four models and five seeds
whose primary estimand is learning-rate invariant and whose selector is fixed by
equal-weight marginalization. Fourth, we make resolving power a preregistered,
executable gate: the design states the effect it can detect, measures the seed
variance it depends on, and refuses the full panel until a pilot confirms it.
Fifth, we provide a layered artifact contract separating design freeze,
development-based checkpoint selection, retrospective evaluation, and a later
prospective confirmation.

